\documentclass[12pt,a4paper]{article}

\usepackage[T2A]{fontenc}
\usepackage[utf8]{inputenc}
\usepackage[russian]{babel}
\usepackage{amsmath, amssymb, amsfonts}
\usepackage{booktabs}
\usepackage{geometry}
\geometry{left=2.5cm, right=2.5cm, top=2cm, bottom=2cm}

\title{Лабораторная работа №3 (МОИ)\\
\large Формирование распределений случайных точек и направлений}
\author{}
\date{}

\begin{document}
\maketitle

\paragraph{Тема.}
Формирование заданных распределений случайных точек и направлений в пространстве.

\paragraph{Цель.}
Реализовать и обосновать корректность четырёх схем сэмплирования: равномерно по треугольнику и диску, равномерно по единичной сфере и по \textbf{косинусно взвешенной} полусфере относительно нормали.

Ниже приведены краткие доказательства в духе учебных пособий по глобальной освещённости; формулировки сознательно отличаются от типовых шаблонов отчётов, логика та же.

\section{Равномерное распределение внутри треугольника}

Пусть треугольник задан вершинами $V_1, V_2, V_3$. Любая точка внутри него однозначно записывается как
\[
P = V_1 + r_1(V_2 - V_1) + r_2(V_3 - V_1), \quad r_1 \ge 0,\; r_2 \ge 0,\; r_1 + r_2 \le 1.
\]

Пару $(r_1, r_2)$ удобно рассматривать как точку в квадрате $[0,1]^2$: независимо берутся $U_1, U_2 \sim \mathcal{U}(0,1)$. Подмножество, соответствующее треугольнику (условие $r_1+r_2 \le 1$), занимает половину площади квадрата. Чтобы не терять равномерность по площади треугольника, пары из «верхнего» треугольника ($r_1+r_2>1$) \textbf{отражают} относительно точки $(1,1)$:
\[
(r_1, r_2) \leftarrow (1-r_1,\, 1-r_2).
\]

Отображение биективно переводит верхний треугольник на нижний, сохраняя меру Лебега (якобиан по модулю равен~1). После отражения $(r_1,r_2)$ равномерно распределена в стандартном симплексе $\{r_1,r_2 \ge 0,\; r_1+r_2\le 1\}$, а значит барицентрические веса
\[
a = 1-r_1-r_2,\quad b = r_1,\quad c = r_2
\]
дают \textbf{равномерное по площади} распределение $P = aV_1 + bV_2 + cV_3$.

\paragraph{Проверка в программе.}
Для выборки проверяется $a,b,c \ge 0$, $a+b+c=1$ и принадлежность точек треугольнику (через вычисление барицентрических координат относительно рёбер).

\paragraph{Замечание про гистограмму по координате $b$.}
У равномерного по площади треугольника маргинальная плотность $b$ на $[0,1]$ равна $f_B(b) = 2(1-b)$ --- она \textbf{убывает} к правому краю. Это согласуется с визуализацией «плотность vs $b$».

\section{Равномерное распределение внутри круга радиуса $R_c$}

В плоскости с ортонормированным базисом $(\hat t, \hat b)$ точка
\[
P = C + \rho\cos\varphi\,\hat t + \rho\sin\varphi\,\hat b.
\]

Равномерность по \textbf{площади} диска достигается, если совместная плотность в полярных координатах удовлетворяет
\[
f_{\rho,\varphi}(\rho,\varphi) = \frac{\rho}{\pi R_c^2}, \quad \rho \in [0,R_c],\; \varphi \in [0,2\pi).
\]

Берутся независимые $U_1, U_2 \sim \mathcal{U}(0,1)$ и задаются
\[
\rho = R_c\sqrt{U_1}, \quad \varphi = 2\pi U_2.
\]

Тогда $F_\rho(r) = P(\rho \le r) = P(U_1 \le r^2/R_c^2) = r^2/R_c^2$, откуда $f_\rho(r) = 2r/R_c^2$; с учётом якобиана $J = \rho$ для перехода к плоским координатам получается равномерность по диску. Вектор нормали $N$ задаёт ориентацию плоскости; базис $(\hat t,\hat b,\hat n)$ строится ортогонализацией (как в \texttt{utils.py}).

\paragraph{Проверка.}
Все точки лежат в круге $\rho \le R_c$; гистограмма $\rho$ согласуется с $f_\rho(r) = 2r/R_c^2$.

\section{Равномерное распределение направлений на единичной сфере $S^2$}

Пусть ось $z$ произвольна. Площадь сферического пояса между уровнями $z$ и $z+\mathrm{d}z$ пропорциональна только $\mathrm{d}z$ (теорема Архимеда). Значит, для равномерности по поверхности сферы координата $z$ должна быть \textbf{равномерной} на $[-1,1]$:
\[
z = 2U_1 - 1, \quad U_1 \sim \mathcal{U}(0,1), \quad \varphi = 2\pi U_2.
\]

Долгота $\varphi$ равномерна на полном круге; радиус в экваториальной плоскости $r_{xy} = \sqrt{1-z^2}$. Тогда
\[
\boldsymbol{\omega} = (r_{xy}\cos\varphi,\; r_{xy}\sin\varphi,\; z)
\]
--- единичный вектор, и распределение \textbf{равномерно по мере на $S^2$}.

\paragraph{Проверка.}
$\|\boldsymbol{\omega}\| = 1$; гистограмма $z$ согласуется с плотностью $f_Z(z) = \tfrac{1}{2}$ на $[-1,1]$.

\section{Косинусное распределение на полусфере (метод Малли)}

Нужно сэмплировать единичные направления $\boldsymbol{\omega}$ в полупространстве $\boldsymbol{\omega}\cdot \hat n \ge 0$ с плотностью
\[
p(\boldsymbol{\omega}) = \frac{1}{\pi}\max(0,\; \boldsymbol{\omega}\cdot \hat n),
\]
что соответствует диффузной (ламбертовой) модели в Монте-Карло.

В локальном ортонормированном базисе $(\hat t,\hat b,\hat n)$ пишем $\boldsymbol{\omega} = x\hat t + y\hat b + z\hat n$ с $z = \cos\theta \in [0,1]$. Стандартная конструкция (метод Малли):
\[
U_1, U_2 \sim \mathcal{U}(0,1),\quad
z = \sqrt{U_1},\quad
\sin\theta = \sqrt{1-U_1},\quad
\varphi = 2\pi U_2,
\]
\[
x = \sin\theta\cos\varphi,\quad y = \sin\theta\sin\varphi.
\]

Маргинальная плотность \textbf{косинуса угла} $C = \cos\theta = z$ на $[0,1]$ получается из $P(Z \le c) = P(\sqrt{U_1} \le c) = c^2$, откуда
\[
f_C(c) = \frac{\mathrm{d}}{\mathrm{d}c} c^2 = 2c,\quad c \in [0,1].
\]

То есть плотность по $c = \cos\theta$ \textbf{монотонно возрастает} от $0$ до $2$ --- больше направлений укладывается около нормали, что и требуется косинусному взвешиванию.

\subsection{Почему «должно убывать» легко перепутать}

\begin{enumerate}
  \item \textbf{Убывающая кривая в отчёте по треугольнику} относится к маргинали \textbf{$b$} (см. §1): там $f_B(b)=2(1-b)$ действительно убывает. Это другая задача.

  \item Для \textbf{равномерного} распределения по \textbf{полусфере} (не косинусного!) маргинальная плотность $C = \cos\theta$ была бы \textbf{равномерной} на $[0,1]$ ($f_C=1$) --- «плоская» гистограмма, а не косинусная $2c$.

  \item Если строить гистограмму не по $c=\cos\theta$, а, например, по \textbf{углу} $\theta \in [0,\pi/2]$, то совместная плотность $p(\theta,\varphi)\propto \cos\theta\sin\theta$ даёт другую форму кривой по $\theta$ --- её максимум не на $\theta=0$. Путаница осей и величин часто и даёт «не ту» монотонность.
\end{enumerate}

\paragraph{Вывод.}
В корректной постановке \textbf{косинусного} сэмплирования плотность по \textbf{$\cos\theta$} задаётся $f(c)=2c$ и \textbf{возрастает} с $c$. Убывающая плотность в том же виде графика без уточнения оси обычно означает ошибку в постановке опыта или подмену распределения.

\section{Сводная таблица проверок}

\begin{center}
\begin{tabular}{@{}lll@{}}
\toprule
Объект & Ключевой факт & Эмпирическая проверка \\
\midrule
Треугольник & отражение в симплексе & $a,b,c\ge 0$, сумма 1, точки внутри $\triangle$ \\
Диск & $\rho = R_c\sqrt{U}$ & гистограмма $\rho$ согласуется с $2r/R_c^2$ \\
Сфера $S^2$ & $z$ равномерен на $[-1,1]$ & гистограмма $z$ согласуется с $1/2$ \\
Косинус-полусфера & Малли: $z=\sqrt{U_1}$ & гистограмма $\cos\theta$ согласуется с $2c$ на $[0,1]$ \\
\bottomrule
\end{tabular}
\end{center}

\section*{Литература}
Д.\,Д.~Жданов, И.\,С.~Потемин, А.\,Д.~Жданов, \textit{Методы расчёта глобальной освещённости} (ИТМО, 2023).

\paragraph{Примечание.}
При необходимости добавьте скриншоты интерфейса \texttt{labmoi3} и сравнение гистограмм с теоретическими кривыми.

\end{document}
